```latex
\documentclass[11pt]{article}
\usepackage{amsmath,amssymb,amsthm}
\usepackage[margin=1in]{geometry}
\usepackage{hyperref}

\newcommand{\M}{\mathcal M}
\newcommand{\C}{\mathbb C}
\newcommand{\D}{\Delta}
\newcommand{\barM}{\overline{\M}}

\newtheorem{theorem}{Theorem}[section]
\newtheorem{proposition}[theorem]{Proposition}
\newtheorem{lemma}[theorem]{Lemma}
\newtheorem{corollary}[theorem]{Corollary}
\newtheorem{remark}[theorem]{Remark}
\newtheorem{definition}[theorem]{Definition}

\title{On the proposed non-variation of a point factor}
\author{}
\date{}

\begin{document}
\maketitle

\section{The obstruction}

The stated conclusion does not follow from the listed hypotheses in the
local-submanifold sense allowed by the footnote.  The boundary theorems for
totally geodesic subvarieties control the central fibre \(V_0\), but the
hypotheses for \(t\neq 0\) are only smooth and combinatorial.  In particular,
they do not prevent a point factor from moving holomorphically.

More precisely, the boundary theorem for totally geodesic subvarieties says
that an irreducible component of the intersection of the closure of a totally
geodesic subvariety with a Deligne--Mumford stratum is again totally geodesic
and has a prime decomposition whose projections are locally isometric
\cite[Theorems~1.4 and~1.7]{BW}.  A semisimple formulation is given in
\cite[Theorem~1.3]{AW}.  These theorems apply to \(V_0\), but they say nothing
about the arbitrary deformations \(V_t\), which are not assumed to be totally
geodesic.

We first record the relevant boundary theorem.

\begin{theorem}[Boundary semisimplicity]
Let \(N\subset \M_{g,n}\) be totally geodesic, and let
\(\mathring{\partial}N\) be an irreducible component of
\(\partial N\cap \Delta_G\).  Then \(\mathring{\partial}N\) is totally
geodesic in \(\Delta_G\).  Moreover, its product cover decomposes into prime
factors, and the projection from each positive-dimensional prime factor to
every moduli-space factor occurring in it is locally injective, in the orbifold
sense.
\end{theorem}

\begin{proof}
This is Theorems~1.4 and~1.7 in \cite{BW}; equivalently, it is the semisimple
boundary theorem in \cite{AW}.
\end{proof}

Applying this to \(V_0\), the hypothesis
\(\pi_i(\partial V_0)=\{Y_0\}\) says that the \(i\)-th factor occurs as a
split constant factor in the semisimple decomposition.  Thus, locally in the
product cover, one has a model of the form
\[
\partial V_0 \simeq C_0\times \{Y_0\},
\]
possibly with further singleton factors.  The theorem does not give a discrete
characterization of \(Y_0\), and it gives no information about \(Y_t\) for
\(t\neq 0\).

\section{A local deformation moving the point factor}

We now give the standard boundary-coordinate model showing that the remaining
hypotheses allow \(Y_t\) to move.

Choose a regular boundary point
\[
p\in \partial V_0
\]
at which the closure of \(V_0\) is smooth and transverse to the relevant node
directions.  Let \(r\) be the number of edges of \(\Gamma\).  An orbifold
plumbing chart of \(\barM_g\) near \(p\) has coordinates
\[
(s_1,\ldots,s_r,x,y,z),
\]
where
\[
\Delta_\Gamma=\{s_1=\cdots=s_r=0\},
\]
\(x\) is a coordinate along the one-dimensional varying part of
\(\partial V_0\), \(y\) denotes coordinates on the factor
\(\M_{g_i,n_i}\), and \(z\) denotes the remaining factor coordinates.

After shrinking the chart and making a holomorphic change of coordinates, the
germ of \(V_0\) may be written as
\[
V_{0,\mathrm{loc}}
=
\left\{
(x,u):
s_j=u^{m_j}\quad(1\leq j\leq r),\quad
y=0,\quad z=0
\right\},
\qquad u\neq 0,
\]
for suitable positive integers \(m_j\).  Here \(u\) is the common smoothing
parameter and \(y=0\) corresponds to the point \(Y_0\).  This is simply the
holomorphic implicit function theorem applied to the smooth germ of
\(\overline V_0\) along the smooth curve \(\partial V_0\).

Fix a nonzero tangent vector \(a\in T_{Y_0}\M_{g_i,n_i}\).  In a local affine
coordinate on that factor, define, for sufficiently small \(t\),
\[
V_{t,\mathrm{loc}}
=
\left\{
(x,u):
s_j=u^{m_j}\quad(1\leq j\leq r),\quad
y=t a,\quad z=0
\right\},
\qquad u\neq 0.
\]
Then:

\begin{enumerate}
\item each \(V_{t,\mathrm{loc}}\) is a smooth complex surface;
\item the total space
\[
\mathcal V_{\mathrm{loc}}
=
\{(x,u,t):s_j=u^{m_j},\ y=ta,\ z=0,\ u\neq 0\}
\]
is a complex submanifold of \((\M_g\cap U)\times \D\);
\item the family is smoothly trivial, since \((x,u,t)\) are global coordinates
on it;
\item
\[
\overline{V_{t,\mathrm{loc}}}\cap \Delta_\Gamma
=
\{(x,0):s_j=0,\ y=ta,\ z=0\};
\]
hence
\[
\pi_i\bigl(\partial V_{t,\mathrm{loc}}\bigr)=\{Y(t)\},
\qquad
Y(t)=Y_0+t\,a.
\]
\end{enumerate}

By choosing the \(x\)- and \(u\)-coordinate domains relatively compact in the
chart, no other Deligne--Mumford strata are encountered in the local closure.
Thus the local versions of hypotheses (i) and (ii) hold.  Nevertheless
\(Y(t)\) is nonconstant whenever \(a\neq 0\).

\begin{remark}
At \(t=0\), the germ is the germ of the totally geodesic surface \(V_0\); in
particular its second fundamental form vanishes.  The construction changes
only the extrinsic location of a split point factor.  No hypothesis in the
question rules this out for \(t\neq0\).
\end{remark}

This shows that, under the complex-submanifold formulation in the footnote,
the proposed assertion is false as stated.

\section{What additional hypothesis would give the conclusion}

The following is a correct rigidity statement under a substantially stronger
assumption.

\begin{proposition}
Suppose that \(V_t\) is totally geodesic for every \(t\in\D\).  Then
\(V_t=V_0\) for all \(t\).  In particular,
\[
\pi_i(\partial V_t)=\{Y_0\}
\]
for every \(t\).
\end{proposition}

\begin{proof}
Every higher-dimensional totally geodesic submanifold of Teichm\"uller space
has algebraic image in moduli space, and there are only finitely many
higher-dimensional totally geodesic subvarieties in a fixed moduli space
\(\M_g\); see \cite{AW}.  Let
\[
\mathscr F_g
=
\{W\subset\M_g:\dim_{\C}W=2,\ W\text{ is totally geodesic}\}.
\]
This set is finite.

A holomorphic, smoothly trivial family gives a continuous assignment
\[
t\longmapsto V_t
\]
in the space of compactified cycles, or equivalently in the relevant Chow or
Douady space.  Since every \(V_t\) lies in the finite, hence discrete, set
\(\mathscr F_g\), and since \(\D\) is connected, this map is constant.
Therefore \(V_t=V_0\) for every \(t\).  Taking closures and intersecting with
\(\Delta_\Gamma\) gives
\[
\partial V_t=\partial V_0,
\]
and hence
\[
\pi_i(\partial V_t)=\pi_i(\partial V_0)=\{Y_0\}.
\]
\end{proof}

\section{Conclusion}

The boundary semisimplicity theorem for \(V_0\) permits a decomposition
\[
\partial V_0=C_0\times\{Y_0\},
\]
but it does not impose discreteness on the point \(Y_0\).  A deformation of
the form
\[
C_t\times\{Y(t)\}
\]
is compatible with the stated smooth-triviality and incidence conditions.
Consequently, the proposed conclusion requires an additional rigidity
hypothesis---for example, that every \(V_t\) is totally geodesic, or more
generally that the family lies in a discrete set of admissible subvarieties.

\begin{thebibliography}{9}

\bibitem{BW}
B.~Wright,
\emph{The boundary of a totally geodesic subvariety of moduli space},
available at
\href{https://arxiv.org/html/2412.06765v1}{\texttt{arxiv.org/html/2412.06765v1}}.

\bibitem{AW}
F.~Arana-Herrera and A.~Wright,
\emph{The geometry of totally geodesic subvarieties of moduli spaces of
Riemann surfaces},
available at
\href{https://arxiv.org/html/2412.16330v2}{\texttt{arxiv.org/html/2412.16330v2}}.

\end{thebibliography}

\end{document}
```